%% kapa: exam=2 round=26 year=2015 period=1 subject=appraisal_practice subject_ko=감정평가실무 seq=4 number=4 points=10
%% source: Q-Net 감정평가사 기출문제 (cst00309) artlSeq=5149601 file=extracted/1008/5149601/2015년 제26회 감정평가사 2차시험 기출문제/1교시(감정평가실무).pdf
%% conversion: kordoc 4.24.51 -> tools/gen2.py (auto); review=proofread-gpt-sol; flags=figure
\begin{longquestion}{【문제 4】}
감정평가사 김공정씨는 다음 물건에 대하여 ○○지방법원으로부터 경매 목적의 감정평가를 의뢰받았다. 기준시점을 2015.09.19로 하여 관련법규 및 이론을 참작하고 주어진 자료를 활용하여 감정평가하시오. (10점) \par
<자료 1> 법원감정평가 명령서 내용 요약 및 평가대상 \par
1. 기호(1): S시 S구 S동 1210번지, 대, 200 m\textsuperscript{2}, 제2종일반주거지역 \par
2. 기호(가): S시 S구 S동 1210번지 지상 철근콘크리트조 및 벽돌조 슬래브지붕 2층 주택(사용승인일: 2009.12.05, 완공일: 2008.05.05) \par
1층 : 철근콘크리트조 단독주택 100 m\textsuperscript{2} \par
2층 : 벽돌조 단독주택 12 m\textsuperscript{2} (2012.02.03 증축) \par
3. 현장조사사항 : 기호(1) 토지 지상에는 <자료 2> 현황도와 같이 법원의 제시목록 외 기호㉠, ㉡이 소재함. 제시목록뿐 아니라 등기사항전부증명서 및 대장에도 등재되어 있지 아니하여 소유권에 대한 재확인이 필요함 \par
4. 유의사항 : 제시 외 건물이 있는 경우에는 반드시 그 가액을 평가하고, 제시 외 건물이 경매대상에서 제외되어 그 대지가 소유권의 행사를 제한 받는 경우에는 그 제한을 반영하여 평가함 \par
<자료 2> 현황도 \par
\qfigure[0.78\linewidth]{figures/04-1.png}

<자료 3> 건물평가자료 \par
{\small\centering
\resizebox{\ifdim\width>\linewidth\linewidth\else\width\fi}{!}{%
\begin{tabular}{|c|c|c|c|c|c|}
\hline
\makecell{구분} & \makecell{구조} & \makecell{이용상황} & \makecell{재조달원가\\(원/m\textsuperscript{2})} & \makecell{면적\\(m\textsuperscript{2})} & \makecell{적용단가\\(원/m\textsuperscript{2})} \\
\hline
\makecell{기호(가)기존} & \makecell{철근콘크리트조\\슬래브지붕} & \makecell{주택(방1, 거실,\\주방, 화장실1)} & \makecell{750,000} &  &  \\
\hline
\makecell{기호(가)증축} & \makecell{벽돌조\\슬래브지붕} & \makecell{방1} & \makecell{600,000} &  &  \\
\hline
\makecell{제시 외\\건물 기호㉠} & \makecell{경량철골조\\판넬지붕} & \makecell{보일러실} & \makecell{-} & \makecell{4} & \makecell{100,000} \\
\hline
\makecell{제시 외\\건물 기호㉡} & \makecell{벽돌조\\슬래브지붕} & \makecell{주택(방1, 주방1,\\화장실1)} & \makecell{600,000} & \makecell{48} &  \\
\hline
\end{tabular}}\par}

※ 철근콘크리트조 내용년수 50년, 벽돌조 내용년수 45년, 잔가율 0 \% \par
※ 제시 외 건물 기호㉠은 신축연도가 불명확하여 관찰감가를 병용하여 적용단가를 산정하였으며, 면적은 실측면적임 \par
※ 제시 외 건물 기호㉡은 잔존내용년수가 20년으로 추정됨 \par
<자료 4> 기타자료 \par
1. 제시 외 건물이 토지에 미치는 영향을 고려하지 아니하고 공시지가 기준으로 평가한 금액은 6,530,000 원/m\textsuperscript{2}임. 제시 외 건물이 토지에 미치는 영향이 있다고 판단될 경우에는 아래사항을 감안하여 평가하기 바람 \par
{\small\centering
\resizebox{\ifdim\width>\linewidth\linewidth\else\width\fi}{!}{%
\begin{tabular}{|c|c|}
\hline
 & \makecell{전체 토지에 미치는 영향} \\
\hline
\makecell{제시 외 건물 기호㉠} & \makecell{1 \%} \\
\hline
\makecell{제시 외 건물 기호㉡} & \makecell{12 \%} \\
\hline
\end{tabular}}\par}

\end{longquestion}
